% !PRINTED 405
\weberpage{405}
%<<CONT

and in the two exceptional cases $D=-4$, $D=-3$ the last
expression is to be divided by 2 or by 3.

\gesperrt{Here $\mathfrak{m}$ runs through the totality of the ideals of a
class (with respect to $[Q]$), and one sees that this limiting
value does not depend on this particular class.}

It is sometimes expedient to exclude from the sum on the left-hand
side of formula (12) not merely those ideals of $\mathfrak{m}$
which have a divisor in common with $Q$, but also those which are
not coprime to $\varDelta$. Formula (12) holds unchanged under
this assumption if by the $r$ of the right-hand side one
understands not only the prime divisors of $Q$, but all the prime
divisors of $D$.

For if $\mathfrak{r}$ is a prime ideal which divides $D$ but not
$Q$, then $N(\mathfrak{r}) = r_1$ is a natural prime which divides
$D$ but not $Q$, and $N(\mathfrak{r}\mathfrak{m}) = r_1\,N(\mathfrak{m})$. If, then,
one further removes from the sum on the left-hand side of $y$ the % [sic] The German source prints “von y”; syntactically a reference to the left-hand sum is expected.
terms corresponding to the $\mathfrak{r}\mathfrak{m}$, then on the
right the factor

$\Bigl(1 - \frac{1}{r_1}\Bigr)$ is added.

5. \gesperrt{The limiting value in formula (12) depends only on
the discriminant $D$, not on the particular class $M$,
and is always different from zero.}

\begin{center}
\S\,112. \gesperrt{Class Number.}
\end{center}

We now assume a function $F(z)$, which moreover need only be
defined for integer values of $z$, and of which we suppose that
the infinite series

\weq{1}{\overset{\mathfrak{a}}{\Sigma} F(N\mathfrak{a}),}
in which $\mathfrak{a}$ runs through the totality of the ideals of
the field $\Omega$ (or only a part of them), is \gesperrt{unconditionally
convergent}. When it is a question of orders $[Q]$, we wish to
\gesperrt{exclude from the $\mathfrak{a}$ all those which are not relatively
prime to $Q$.}

Any natural number $m$ coprime to $Q$ then occurs among the
$N(\mathfrak{a})$, by § 93, 4.

\[ \overset{n}{\Sigma} (\varDelta, n) \]
times, if $n$ runs through the divisors of $m$, and we obtain
the formula:

\weq{2}{\sum F(N\mathfrak{a}) = \overset{m}{\Sigma} \overset{n}{\Sigma} (\varDelta, n)\,F(m),}
%CONT>>
